\documentclass[11pt,a4paper]{article}

\usepackage[margin=2.5cm]{geometry}

\title{Hybrid Quantum Computing with Haskell and QPP V3}
\author{}
\date{\today}

\begin{document}

\maketitle

\tableofcontents

%%%%%%%%%%%%%%%%%%%%%%%%%%%%%%%%%%%%%%%%%%%%%%%%%%%%%%%%%%%%%
\part{Quantum Programming in QPP}

\section{Introduction}

\subsection{The Goals of the Course}

\subsection{Quantum Computing as Representation}

\subsection{The Quantum Programming Playground}

\subsection{Running QPP}

\subsection{States, Operators, and Programs}


\section{Enough Haskell to Begin}

\subsection{Expressions, Types, and Functions}

\subsection{Complex Numbers}

\subsection{Lists and Comprehensions}

\subsection{Partial Application}

\subsection{Algebraic Data Types}

\subsection{Pattern Matching}

\subsection{Structural Recursion}

\subsection{Exercises}


\section{States and Amplitudes}

\subsection{A Qubit is a Vector}

\subsection{Pure States and Normalization}

\subsection{The Born Rule}

\subsection{Inner Products}

\subsection{Change of Basis}

\subsection{Global and Relative Phase}

\subsection{Exercises}


\section{Operators}

\subsection{Linear Operators}

\subsection{Unitary Evolution}

\subsection{The Pauli Operators}

\subsection{The Hadamard Operator}

\subsection{Rotation Operators}

\subsection{Tensor Products}

\subsection{Controlled Operators}

\subsection{Adjoints}

\subsection{Exercises}


\section{Operators as Syntax Trees}

\subsection{The \texttt{QOp} Datatype}

\subsection{Constructors as Syntax}

\subsection{Syntax versus Semantics}

\subsection{The Matrix Backend}

\subsection{Operator Composition}

\subsection{Tensor Composition}

\subsection{Structural Recursion on Operator Syntax}

\subsection{Analysis of Operator Expressions}

\subsection{Exercises}


\section{Many-Qubit Systems}

\subsection{Tensor-Product Spaces}

\subsection{Product States}

\subsection{Entangled States}

\subsection{Bell States}

\subsection{Bell-State Preparation}

\subsection{Bell-Basis Measurement}

\subsection{Exercises}


\section{Programs with Measurement}

\subsection{The \texttt{Program} Datatype}

\subsection{Unitary and Non-Unitary Operations}

\subsection{Measurement as an Effect}

\subsection{Randomness and Random Streams}

\subsection{State Collapse}

\subsection{Classical Control}

\subsection{Adaptive Quantum Programs}

\subsection{Sampling and Reproducibility}

\subsection{Exercises}


%%%%%%%%%%%%%%%%%%%%%%%%%%%%%%%%%%%%%%%%%%%%%%%%%%%%%%%%%%%%%
\part{Hybrid Quantum--Classical Computation}

\section{The Input Problem}

\subsection{State Preparation}

\subsection{Preparing Structured States}

\subsection{Entangled-State Preparation}

\subsection{Complexity of State Preparation}

\subsection{Exercises}


\section{The Output Problem}

\subsection{Measurement as Information Extraction}

\subsection{Amplitude Amplification}

\subsection{Grover Search}

\subsection{Output Complexity}

\subsection{Exercises}


\section{The Simulation Problem}

\subsection{Why Simulation is Difficult}

\subsection{Exponential State Spaces}

\subsection{Restricted Quantum Computation}

\subsection{The Clifford Group}

\subsection{Clifford versus Non-Clifford Computation}

\subsection{Exercises}


%%%%%%%%%%%%%%%%%%%%%%%%%%%%%%%%%%%%%%%%%%%%%%%%%%%%%%%%%%%%%
\part{Quantum Operators as Algebraic Objects}

\section{The Pauli Algebra}

\subsection{The Pauli Group}

\subsection{Multiplication Rules}

\subsection{Commutation and Anticommutation}

\subsection{Multi-Qubit Pauli Operators}

\subsection{Generating Sets}

\subsection{Operator Bases}

\subsection{Normal Forms}

\subsection{Exercises}


\section{Structural Transformations}

\subsection{Traversing Operator Trees}

\subsection{Analysis versus Transformation}

\subsection{Simplification Rules}

\subsection{Rewrite Systems}

\subsection{Canonical Forms}

\subsection{Validation and Consistency Checking}

\subsection{Correctness of Rewrites}

\subsection{Exercises}


\section{Interpreters and Semantics}

\subsection{Syntax and Interpretation}

\subsection{Matrix Semantics}

\subsection{State-Vector Semantics}

\subsection{Symbolic Semantics}

\subsection{Resource-Analysis Semantics}

\subsection{Comparing Interpretations}

\subsection{Exercises}


\section{Why Representations Matter}

\subsection{Representation-Dependent Complexity}

\subsection{State Vectors versus Symbolic Forms}

\subsection{Choosing a Representation}

\subsection{Representation Trade-Offs}


%%%%%%%%%%%%%%%%%%%%%%%%%%%%%%%%%%%%%%%%%%%%%%%%%%%%%%%%%%%%%
\part{Alternative Quantum Representations}

\section{The Stabilizer Formalism}

\subsection{Stabilizer Groups}

\subsection{Stabilizer States}

\subsection{Generators}

\subsection{The Gottesman--Knill Theorem}

\subsection{Examples}

\subsection{Exercises}


\section{Tableau Representations}

\subsection{From Stabilizers to Tableaux}

\subsection{Pauli Strings}

\subsection{Tableau Data Structures}

\subsection{Representing Stabilizer States}

\subsection{Hadamard Updates}

\subsection{Phase-Gate Updates}

\subsection{CNOT Updates}

\subsection{Measurement}

\subsection{Examples}

\subsection{Exercises}


\section{Diagrammatic Representations}

\subsection{Motivation}

\subsection{Circuit Diagrams}

\subsection{String Diagrams}

\subsection{ZX Diagrams}

\subsection{Z-Spiders}

\subsection{X-Spiders}

\subsection{Hadamard Nodes}

\subsection{Diagram Semantics}

\subsection{Examples}


\section{Graph Rewriting}

\subsection{Graphs as Syntax}

\subsection{Rewrite Rules}

\subsection{Spider Fusion}

\subsection{Identity Removal}

\subsection{Hadamard Simplifications}

\subsection{Simplification Strategies}

\subsection{Termination and Confluence}

\subsection{Circuit Extraction}

\subsection{Examples}


\section{Tensor Representations}

\subsection{Operators as Tensors}

\subsection{Tensor Networks}

\subsection{Contraction}

\subsection{Complexity Considerations}

\subsection{Simulation Strategies}

\subsection{Examples}

\subsection{Exercises}


%%%%%%%%%%%%%%%%%%%%%%%%%%%%%%%%%%%%%%%%%%%%%%%%%%%%%%%%%%%%%
\part{Projects and Research Directions}

\section{Project Track: Symbolic Operator Algebra}

\subsection{Operator Simplifiers}

\subsection{Canonical Forms}

\subsection{Commutation Analysis}

\subsection{Verification Tools}


\section{Project Track: Stabilizer Computation}

\subsection{Stabilizer Simulators}

\subsection{Tableau Simulators}

\subsection{Clifford Evolution}

\subsection{Measurement}

\subsection{Error-Correction Applications}


\section{Project Track: ZX Calculus}

\subsection{ZX Graph Representations}

\subsection{Rewrite Engines}

\subsection{Simplification Strategies}

\subsection{Circuit Extraction}


\section{Project Track: Tensor Methods}

\subsection{Tensor-Network Simulators}

\subsection{Contraction Heuristics}

\subsection{Backend Implementations}

\subsection{Performance Comparisons}


\section{Research Directions}

\subsection{Quantum Compilation}

\subsection{Quantum Error Correction}

\subsection{Program Verification}

\subsection{Equivalence Checking}

\subsection{Diagrammatic Reasoning}

\subsection{Quantum Software Architecture}

\subsection{Representation Theory in Quantum Software}


%%%%%%%%%%%%%%%%%%%%%%%%%%%%%%%%%%%%%%%%%%%%%%%%%%%%%%%%%%%%%
\appendix

\part{Optional Topics}

\section{The Bloch Sphere}

\subsection{Global Phase and Projective States}

\subsection{Bloch Coordinates}

\subsection{Pauli Expectation Values}

\subsection{Rotations of the Bloch Sphere}

\subsection{Exercises}


\section{Interference}

\subsection{Hadamard Interference}

\subsection{Phase and Amplitudes}

\subsection{The Mach--Zehnder Interferometer}

\subsection{Interference Fringes}

\subsection{Exercises}


\section{Which-Path Information}

\subsection{Entanglement with a Path Register}

\subsection{Loss of Interference}

\subsection{Decoherence}

\subsection{Quantum Erasure}

\subsection{Exercises}

\end{document}